% 张三慧《学习辅导与习题解答》扫描 S 卷1第7章“思考题选答”原编号。
% 仅辅导书来源已视觉核验，V1主教材原题逐图交叉核验仍待；全部exam_ready=false。
\section*{第7章\quad 波动：辅导书已列思考题独立解析}
\paragraph{思考题 7.6}
可取理想弦驻波 \(y(x,t)=2A\sin(kx)\cos\omega t\)。动能密度 \(u_K=\frac12\mu(\partial_ty)^2=2\mu A^2\omega^2\sin^2kx\sin^2\omega t\)，弹性势能密度 \(u_U=\frac12T_s(\partial_xy)^2=2T_sA^2k^2\cos^2kx\cos^2\omega t\)。若图中 \(t=0\) 为最大形变，\(t=T/4\) 时弦各点通过平衡位置、势能为零而动能最大；\(t=T/2\) 时各点再到极端位置、动能为零而势能最大。节点总不动、局部动能零但弹性变形常最大；波腹位移幅值最大，极端形变时局部斜率为零。驻波的能量在空间和时间交替分布，不能简单把每个质元的“弹性势能”都定义成只依赖其自身位移。

\paragraph{思考题 7.7}
两端有效固定的理想弦基频为 \(\boxed{f_1=\frac{1}{2L}\sqrt{T/\mu}}\)，其中 \(L\) 为有效振动长度、\(T\) 为弦张力、\(\mu\) 为线密度。调弦轴改变张力 \(T\)；手指按弦改变有效长度 \(L\)，使 \(f_1\) 及相应谐波频率变化。真正的二胡发声还与琴码、共鸣箱、弓擦的非线性激励和弦刚度有关。

\paragraph{思考题 7.14}
在空气整体匀速漂移且警笛、听众相对于地面都静止的近似中，声波向南相对地面的传播速度由 \(c\) 变为 \(c+u\)；但源每周期 \(T\) 发一个同相波峰，波长同时变成 \(\lambda'=(c+u)T\)，接收频率 \(f'=(c+u)/\lambda'=1/T=f\)。因此\(\boxed{f'=f}\)，只是传播时延与波长改变。若风速空间不均、声源或观察者相对空气运动或介质压缩，则另需讨论。

\paragraph{思考题 7.17}
单色相速度 \(v_{\rm ph}=c/n\) 描述固定相位传播；窄带波包的包络群速度为 \(\boxed{v_g=d\omega/dk=c/[n+\omega(dn/d\omega)]}\)，可等价写成 \(v_g=v_{\rm ph}-\lambda\,dv_{\rm ph}/d\lambda\)。若色散明显，群速度不等于相速度，辅导书比较约 \(1.83\times10^8\) 与 \(1.70\times10^8\,\mathrm{m/s}\) 属不同速度定义。含吸收和反常色散的脉冲前沿及信息传播还需更严谨的因果描述，不能机械断言任意群速度都等于信息速度。
